% !TEX root = ../ai-agent-note-1.tex
\chapter{Agent,工作流和Harness}

\section{谁决定下一步：工作流与 Agent}

\[Agent=LLM+\text{上下文}+\text{工具}\]
LLM 是大脑，负责做决策；上下文是眼睛，决定模型这一步能看到什么；工具是手脚，用来查询或者改变外部世界.\\

\begin{tcolorbox}[colback=second!5, colframe=second, title=Agent\&工作流]
  Context、Tools、Constrain、Verify、Correct
  代码写死步骤，就是工作流(Workflow)；由 LLM 在运行时决定下一步，就是 Agent。
\end{tcolorbox}
两者都是 LLM 加其他组件组成的系统（compound AI system），区别只在流程由谁决定。\\

用代码解释:

\textbf{1.工作流：}控制流写在代码里，LLM 只是某个函数内部的一次调用。
\begin{lstlisting}[style=agentpython, caption={工作流}, label={lst:workflow}]
  def book_flight(req):
    user = verify_identity(req)
    flights = search_flights(req)   # 这里面可以调 LLM
    pay(user, flights[0])           # “付款前不能预订”由代码保证
    return confirm(user)
\end{lstlisting}

\textbf{2.Agent:}控制流变成了模型在运行时吐出来的数据。代码只剩一个循环，每一圈都问模型：“下一步做什么？”
\begin{lstlisting}[style=agentpython, caption={agent}, label={lst:agent-loop}]
  while True:
    action = llm(context)           # 由模型决定下一步
    if action.is_final:
        return action.answer
    obs = execute(action)           # 程序真正执行这个动作
    context.append(action, obs)     # 执行结果进入下一轮上下文
\end{lstlisting}

\section{工具调用怎么发生}
模型从来不执行工具，它只“说”自己想调用哪个工具、传什么参数。真正去执行的是你的程序。\\

\textbf{模型调用的四步流程：}
1.声明工具 2.模型决定调用 3.结果追加上下文 4.模型基于结果回复\\

代码举例：以查天气为例子\\

\textbf{第 一 步：声明工具} 你的程序发出第一次请求，除了消息列表，还要在单独的 tools 字段里告诉模型有哪些工具可用：
\begin{lstlisting}[style=agentpython]
  {
  "messages": [
    {"role": "system", "content": "你是天气助手，需要实时数据时调用工具。"},
    {"role": "user",   "content": "温哥华现在多少度？"}
  ],
  "tools": [{
    "name": "get_weather",
    "description": "查询城市的实时天气",
    "parameters": {"type": "object",
                   "properties": {"city": {"type": "string"}},
                   "required": ["city"]}
  }]
}
\end{lstlisting}
\textbf{第 二 步：模型决定调用:} 模型不回答问题，而是返回一条带 tool\_calls 的 assistant 消息，然后这次请求就结束了。模型没法停在那里等结果。

\begin{lstlisting}[style=agentpython]
{
{"role": "assistant", "content": null,
 "tool_calls": [{"id": "call_1",
                 "function": {"name": "get_weather",
                              "arguments": "{\"city\": \"Vancouver\"}"}}]}
}
\end{lstlisting}
\textbf{第三步:程序执行，把结果追加到上下文。}代码解析 tool\_calls，校验参数，真的去调天气接口，然后往消息列表里追加两条消息：
\begin{lstlisting}[style=agentpython]
{
{"role": "assistant", "tool_calls": [ ...第 ② 步那条，原样放回... ]},
{"role": "tool", "tool_call_id": "call_1", "content": "{\"temp\": 13, \"conditions\": \"晴\"}"}
\end{lstlisting}
\textbf{第四步：模型基于结果回复。} 你的程序发出第二次请求，把完整的消息列表重新发一遍。模型这时才看到结果，回答：“温哥华现在晴，13°C。”

\begin{tcolorbox}[colback=second!5, colframe=second, title=三条规则]
  \begin{enumerate}
    \item API 是无状态的。模型不记得上一次请求里发生过什么，每次请求都要把完整历史重新发送。所以第 1 节 Agent 循环里那句 \texttt{context.append(...)} 不是可有可无，没有它模型就失忆了。
    \item \texttt{tool\_call\_id} 必须配对。一条 assistant 消息里可以同时有好几个工具调用。书里多币种汇总的例子，第一轮就一次发起了 3 次汇率转换。每条 tool 消息要靠 id 说明自己是哪一个调用的结果。
    \item 不要把工具结果伪装成 user 消息。模型会以为是用户又问了一个新问题，多步推理就接不上了。
  \end{enumerate}
\end{tcolorbox}
\section{Harness}

Harness 指的是 Agent 内部、模型之外的所有代码，负责让 Agent 跑起来、并且管住它。第 2 节里“解析 \texttt{tool\_calls}、执行工具、追加消息”的那部分代码，就是 Harness 最基本的一层\\

\begin{tcolorbox}[enhanced, sharp corners, boxrule=0pt, frame hidden,
                  borderline west={3pt}{0pt}{structurecolor},
                  colback=structurecolor!5]
  Agent = Model + Harness\\
  Harness = 上下文管理 + 工具接口 + 约束 + 验证 + 纠正
\end{tcolorbox}

边界要分清：Harness 不包括环境。沙盒里的文件、数据库、用户、外部网站都属于环境（Environment）。Harness 可以创建或隔离环境，但环境本身不归它。

\begin{table}[htbp]
\centering
\caption{Harness 五要素}
\label{tab:harness}
\begin{tabular}{p{2.6cm} p{4.2cm} p{3.6cm} p{4.2cm}}
  \toprule
  要素 & 职责与原则 & 内核里对应什么 & 例子 \\
  \midrule
  \textbf{Context} 上下文   & 信息要充分 & 给进程准备好它需要的数据 & 系统提示词、知识库、历史压缩 \\
  \textbf{Tools} 工具       & 接口要清晰 & 系统调用表 & MCP 工具、代码解释器 \\
  \textbf{Constrain} 约束   & 故障安全，默认关闭 & 权限检查 & 危险工具默认要用户批准 \\
  \textbf{Verify} 验证      & 检查结构化数据，不检查自由文本 & 校验参数和返回值 & linter、类型检查、跑测试 \\
  \textbf{Correct} 纠正     & 静默重试，确认无法恢复之前不暴露中间状态 & 错误处理和恢复 & 重试、熔断、转人工 \\
  \bottomrule
\end{tabular}
\end{table}\newpage

下面是控制循环的简易代码
\begin{lstlisting}[style=agentpython, caption={控制循环}]
  step = 0
  while step < MAX_STEPS:                         # 停止条件，防止死循环
    action = llm(build_context(history))        # Context：决定模型这一轮看到什么
    if action.is_final:
        return action.answer
    if not allowed(action):                     # Constrain：没权限就拦下
        history.append(action, "被拒绝：需要用户批准")
        step += 1
        continue
    obs = execute(action)                       # Tools：真正执行
    if not verify(obs):                         # Verify：检查结果对不对
        obs = correct(action, obs)              # Correct：重试或者换路
    history.append(action, obs)
    step += 1
\end{lstlisting}

\section{agent工程范式}
agent工程重心演变由提示工程到Graph工程

\begin{tcolorbox}[enhanced, sharp corners, boxrule=0pt, frame hidden,
                  borderline west={3pt}{0pt}{structurecolor},
                  colback=structurecolor!5]
\[
  \text{提示工程} \rightarrow \text{上下文工程} \rightarrow \text{Harness 工程}
  \rightarrow \text{Loop 工程} \rightarrow \text{Graph 工程}
\]
\end{tcolorbox}

\begin{table}[htbp]
\centering
\caption{从提示工程到 Graph 工程}
\label{tab:engineering-layers}
\begin{tabular}{p{3cm} p{10cm}}
  \toprule
  层次 & 关注什么 \\
  \midrule
  提示工程     & 一次调用的提示词怎么写 \\
  上下文工程   & 每次调用往窗口里放哪些信息 \\
  Harness 工程 & 模型外面的整套运行与治理：约束、验证、纠正 \\
  Loop 工程    & 让循环跨轮次、跨会话持续运转 \\
  Graph 工程   & 把循环、确定性程序和人工审批组织成一张执行图 \\
  \bottomrule
\end{tabular}
\end{table}

构建有效Agent的三原则：
\begin{enumerate}
  \item 保持简单：先直接调 API，确实需要时再上框架。

  \item 保持透明：把规划、日志和决策轨迹都暴露出来。

  \item 设计好 Agent 与计算机之间的接口（ACI）：站在 Agent 的角度设计工具，用防呆（Poka-yoke）的思路，让错误根本没机会发生。
ACI 不只是写代码，工具描述也是接口的一部分。\\
防呆的核心：与其写提示词提醒模型“请注意不要犯错”，不如改接口，让那个错误写不出来，或者一写出来就得到清楚的报错。
\end{enumerate}

\section{上下文React轨迹}

调用LLM上下文分为五个部分：
\begin{enumerate}
  \item 系统提示词
  \item 工具定义
  \item 用户消息
  \item 模型回复
  \item 工具执行结果
\end{enumerate}
前两项（系统提示词+工具定义）是静态前缀，后三项（用户消息+模型回复+工具执行结果）是随交互不
断增长的动态消息历史。这五个部分共同构成了LLM每次推理时的上下文。\\       
\[\text{每次调用模型的上下文}=\text{静态前缀}+\text{轨迹}\]

\begin{definition}
  React是Reasoning+Acting缩写，每一圈三步：
  \begin{enumerate}
  \item 思考：模型想下一步该做什么
  \item 行动：发出工具调用
  \item 观察：读取工具结果
  \end{enumerate}
\end{definition}

转完一圈，轨迹多出一段，然后进入下一圈，直到模型直接给出答案。这就是第 1 节那个 while True 循环在做的事。
所以为什么要把思考和行动交替进行？
\begin{enumerate}
\item 只思考不行动（纯 CoT）：模型只能凭记忆推理，容易编造事实。
\item 只行动不思考：盲目调用工具，不知道为什么调、调完怎么用。
\item ReAct：每一步的思考都建立在上一步的观察之上，幻觉更少。
\end{enumerate} 

